\section{Search results and the scope of the negative evidence}
\label{sec:search}

This investigation did not infer a nonhalting theorem from finite testing.
Finite prepared-frame tests serve as fast rejection gates for optimization
recipes. A survivor would still require a universal simulation proof and
an exact sufficient blank prefix. In particular, a failed affine proof
with a fuel limit is an unresolved proof attempt unless a concrete
counterexample is separately supplied. A search over a stated family is
not a lower bound on arbitrary machines.

\subsection{Why shortening the return scanner is a joint compiler problem}

The inherited return scanner distinguishes the program counter from the
long zero gap by looking for a sufficiently long zero run. Its chain is
\[
4,6,9,13,19,29,41,58,74,
\]
followed by state 84, which returns right to the dispatcher. An internal
one restarts the scan at state 4. Removing the final detector state is
an attractive one-state reduction, but maximum zero run in the static
source words is not the whole invariant.

As the prior \emph{Beyond 246} packet establishes, an arithmetic phase
temporarily adds a zero next to a trailing zero run. Eighteen decrement
sites with seven trailing zeros then create a run of eight zeros;
both guards yield 36 bad cases for the shortened scanner. The shortest
recorded problematic source word is
\begin{center}\ttfamily 111010010010000000.\end{center}
The issue already occurs before the arithmetic branch is distinguished.
All these sites lie below the shared prefix
\code{11101001001}. A static source-only count would incorrectly accept
the scanner deletion.

The new experiment applied all 238 options of the inherited closed
reassociation/commutation generator to Context246 and tried the shorter
scanner against each changed source program. Every option failed the
early affine screen; none produced a candidate below 246 states. This
does not show that no equivalent tree can support the shorter scanner.
It only rules out acceptance by this checker of the recorded one-step
rewrites under the existing encoding and interface.

\subsection{An exhausted directed layout search with a defined cost gate}

The next search changes the source tree globally at a shared subtree,
interns equal subtrees, reuses matching entries in the original optimized
decoder, and allocates new states for new subtree shapes. This is the
source27 fingerprint-based emitter applied to the Context246 runtime.
We measured its literal tables, not just the number of source DAG nodes.

Let $g(T)$ be the number of working states syntactically reachable from
old-label start 95. Let $b(T)$ be the number of classes of exact one-step
transition bisimulation on that reachable graph. The implemented cost is
\[
 c(T)=\begin{cases}b(T)&g(T)\leq249,\\g(T)&g(T)>249.\end{cases}
\]
Thus the cost gate can exclude a graph with more than 249 states even
if an uncomputed quotient would be small. The shape gate additionally
requires maximum source zero run at most seven and maximum leaf depth
at most 39. These gates are part of the search definition.

\begin{center}
\begin{tabular}{lrr}
\toprule
Quantity & Ceiling 247 & Ceiling 248\\
\midrule
Expanded layouts &364&991\\
Generated rewrite proposals &86,800&237,088\\
Distinct rooted fingerprints encountered &84,700&230,351\\
Unexpanded queue at stop &0&440\\
Smallest charged state count &246&246\\
Termination &Queue exhausted&Time limit\\
\bottomrule
\end{tabular}
\end{center}

The completed 247-gated run contains 56 accepted layouts of cost 246
(including the initial layout) and 308 of cost 247. Its recorded boundary
includes 776 proposals of cost 248. The claim is that the queue was
exhausted for the specified \emph{fingerprint-deduplicated directed global
shared-subtree rewrite algorithm and gates}. It is not an unqualified connected
component under arbitrary inverse rewrites. Global rewriting followed by
interning can coalesce an altered subtree with a previously existing one;
the nominal inverse global rewrite need not undo exactly the previous
operation. The visited-before-gate implementation is nevertheless valid
for this directed search: the gates and the emitted cost are functions
of the rooted interned layout up to fresh-state renumbering.

The 248-gated run stopped after about 361 seconds with a nonempty queue.
No exhaustive claim applies to it. The reproducibility instructions use
its deterministic 991-layout expansion limit to reproduce the functional
record independently of processor speed. Neither cost search substitutes
for the runtime and initialization checks on a surviving candidate.

\subsection{Reassociation with explicit control-target repair}

A broader one-step experiment removed the closedness precondition.
For each binary reassociation, it first required every increment and
both decrement successors to retain their original leaf targets. It
then tried to re-encode each unconditional jump with an existing break
or continue kind leading to its required new target. Ties were resolved
deterministically, preferring the same opcode and then the smaller
operand. The independent complete leaf-graph certificate was the
acceptance gate for the source program.

There were 331 proposed global reassociations. Exactly 255 passed the
source certificate, 46 failed an arithmetic-successor condition, and
30 could not be repaired using the selected existing control encodings.
These counts are for this deterministic repair policy, not all possible
control encodings. Seven source-equivalent layouts had charged cost
246. Five were exactly the already retained layouts; their 8,391
universal binary cases passed. The other two stopped with affine fuel
exhaustion. They are not certified binary machines, and no inference
that they truly diverge follows from the fuel result. No emitted graph
below 246 was found in this one-step family.

The experiment demonstrates why a source compiler certificate and a
binary runtime certificate must remain separate. Equal source subtrees
can occur in new surrounding tape contexts; reusing an optimized literal
entry based only on its fingerprint is a candidate emitter technique,
not an unconditional correctness theorem.

\subsection{Fusion domains and the role of the full invariant}

For a 246-state table, there are $\binom{246}{2}=30,135$ unordered state
pairs. The provider-row family identifies one pair and independently
chooses each of the merged state's two rows from one of its two old
providers, giving four recipes per pair and 120,540 recipes per table.
The changed targets are redirected consistently. The exact literal
family and rejection order are supplied in code.

The synthesized-row family is different. It starts from a pair for
which one quotient row already agrees, fixes that row, and permits the
other row to be any binary write, direction, and successor among the
245 retained working states and halt. That gives 984 choices. There
are 352 row-sharing pairs for the baseline and layouts 0, 236, and 237;
layout 233 has 353 and layout 235 has 351. Across all six tables this
is 2,078,208 recipes per fixture-domain screen.

Three domains must be distinguished. The first tests all 8,391 guards
with finite prepared frames. The second merely removes the six excluded
guards and uses 8,385 frames. The third respects the full nine-register
product invariant at every source site and uses four representative
vectors per feasible guard, for 33,540 complete frames per layout.
Of the 8,385 earlier guard-only frames, 5,327 violate a certified mask
for another register. Their failures therefore cannot refute a merger
that is required to work only on the product invariant.

\begin{center}
\begin{tabular}{p{53mm}rrp{35mm}}
\toprule
Search on all six layouts & Recipes & Frames/layout & Survivors\\
\midrule
Provider rows, six guards excluded &723,240&8,385&0\\
Provider rows, full product domain &723,240&33,540&0\\
Synthesized row, six guards excluded &2,078,208&8,385&0\\
Synthesized row, full product domain &2,078,208&33,540&0\\
\bottomrule
\end{tabular}
\end{center}

These are repeated screens of overlapping finite recipe families,
not counts of that many distinct, globally inequivalent machines.
All final product-domain frames satisfy the independently replayed
mask invariant. They still form a finite representative set. A failed
full-frame case rules out the exact required segment contract on that
frame; it does not assert that a modified blank run reaches the frame.
The baseline full-domain provider family was already searched in prior
work. The changed layouts and the stronger-domain reruns are the new
applications here.

\subsection{The prefix-carry pulse state}

State 255 supplies the extra prefix-carry pulse in Context246. Deleting
it and modifying its incoming row is an apparent 245-state candidate,
but simple versions do not preserve both increment and decrement. The
new runtime experiments enlarged the repair family in two ways.

The first contains 21 specified two-, three-, and four-row families
over a fixed 100-action alphabet. Its assignment count is
\[
 8\cdot100^2+10\cdot100^3+3\cdot100^4
 =310,080,000.
\]
The second allows the incoming replacement and any one other retained
row to range over 980 nonhalting actions, with 489 choices of the
auxiliary row. It tests
\[
 489\cdot980^2=469,635,600
\]
assignments. Neither family has a survivor under its recorded finite
register-operation fixtures. The first family compares exact spare
suffixes. The second permits the number of spare blocks to change,
but requires a nonempty suffix consisting only of value-one blocks,
with all nine primary values exact; its inputs also have only
value-one spares. This is a different, weaker spare interface. The exact
per-family state/row choices, action alphabets, fixture counts, and
bounds are in the runtime README and JSON records. The two families
overlap, so their counts are not added as a unique search population.

This evidence applies to the preserved prefix-carry interface. It
does not show that the pulse state is indispensable after changing
the tape encoding, register boundaries, or source macro interface.
The prior 244-state forward-carry attempt has a stronger but still
local obstruction: its zero and one inputs can reach the same complete
configuration after consuming a value-one spare, although their
required branches differ. A deterministic common return cannot then
recover the lost information. That prior collision argument must not
be generalized to all smaller encodings.

\subsection{Where the remaining opportunity lies}

The most concrete new capability is to feed the full certified context
to the compiler and runtime checker. It supports both a harmless extra
zero-decrement test and source-level removal of an always-zero loop.
The following section states the latter certificate precisely. A
source edit must then be costed after decoder sharing, rather than
valued by its number of removed instructions.

A more radical next design would combine a different arithmetic phase
encoding with the decoder, so that no intermediate tape creates the
extra long zero run that currently prevents scanner shortening. The
two-bit operation-phase protocol would have to retain whether it is
initializing, selecting a register, or executing an operation; moving
that information onto tape is useful only if its reader/writer costs
less than the state saved. This is a proposed direction, not an
implemented smaller machine or a proved saving. The supplied finite
leaf maps, invariant-aware macro oracle, complete-frame contracts, and
startup recognizer give explicit acceptance criteria for pursuing it.
